% ch-a4-07.tex —— 第 4 篇（Emerton–Gee–Hellmann）续片，底本 PDF 252–263 = 印刷页 239–250
\ysrc{239}
现在我们给出与命题 5.3.28 类似的描述，对象是 $\XX_P^{\mathrm{dR},+}$ 的完备化；并用它来定义该形式完备化的若干「紧化」。考虑空间
\begin{equation}\tag{5.3.31}\label{eq:a4-5-3-31}
\XX_{B,P}=\wt{\frakg}\times_{\frakg}\wt{\frakg}_P=\bigcup_{w\in W_P}V_w\subseteq\frakg\times G/B\times G/P,
\end{equation}
其中 $W=\prod_\tau\mathcal{S}_d$ 是 $\mathrm{Res}_{F/\Q_p}\mathrm{GL}_d$ 的 Weil 群，而 $W_P\subseteq W$ 是 $W/W_M$ 中元素的最短长度代表元组成的集合（$W_M\subseteq W$ 是 $\mathrm{Res}_{F/\Q_p}M$ 的 Levi 子群 $\mathrm{Res}_{F/\Q_p}P$ 的 Weyl 群）。对 $w=(w_\tau)_\tau\in W_P$，空间 $V_w$ 定义为 Bruhat 胞腔 $G(1,w)\subseteq G/B\times G/P$ 的原像（回忆 $G=\mathrm{Res}_{F/\Q_p}\mathrm{GL}_d$）。在 $V_w$ 上，$G/\mathrm{Res}_{F/\Q_p}B$ 上的泛旗与 $G/\mathrm{Res}_{F/\Q_p}P$ 上的泛旗之交，给出 $\mathrm{Res}_{F/\Q_p}M/\mathrm{Res}_{F/\Q_p}B_M$ 中一个良定义的旗。给定 $w$，记 $w\lambda$ 为 $w$ 对 $\lambda$ 的置换。此外，我们可以把 $w\lambda$ 看成 $M$ 的一个权。于是对每个 $w\in W$，得到下图
\begin{equation}\tag{5.3.32}\label{eq:a4-5-3-32}
\begin{tikzcd}[row sep=large,column sep=large]
 & V_w/G\arrow[dl]\arrow[dr] & \\
 \wt{\frakg}/G & & \wt{\frakm}/\mathrm{Res}_{F/\Q_p}M,
\end{tikzcd}
\end{equation}
其中按记号滥用，我们以 $\wt{\frakm}$ 记 $\mathrm{Res}_{F/\Q_p}M$ 的 Lie 代数。

\begin{prop}{5.3.33}\label{thm:a4-5-3-33}
函子 $D_{p\mathrm{dR}}$ 诱导下图中的同构
\[
\begin{tikzcd}[row sep=large,column sep=large]
 & (\XX_P^{\mathrm{dR},+})^{\wedge}_{ww_0\lambda}\arrow[dl]\arrow[dr] & \\
 (\XX_d^{\mathrm{dR},+})^{\wedge}_{\lambda} & & (\XX_M^{\mathrm{dR},+})^{\wedge}_{ww_0\lambda},
\end{tikzcd}
\]
（即取 (5.3.32) 的图在幂零自同态构成的闭子空间处的完备化）。这里 $(\XX_P^{\mathrm{dR},+})^{\wedge}_{ww_0\lambda}$ 是 $\XX_P^{\mathrm{dR},+}$ 沿 $(\XX_M^{\mathrm{dR},+})^{\wedge}_{ww_0\lambda}\subseteq(\XX_M^{\mathrm{dR},+})$ 在典范映射 $\XX_P^{\mathrm{dR},+}\to\XX_M^{\mathrm{dR},+}$ 下的原像所作的完备化。
\end{prop}

在图 (5.3.32) 中，有两种方式把 $V_w/G$ 补成 $\wt{\frakg}/G$ 上的真叠（proper stack），而同样地，两种紧化中只有一种允许到 $\wt{\frakm}/\mathrm{Res}_{F/\Q_p}M$ 的态射。最简便的方式是考虑 $V_w$ 在 \eqref{eq:a4-5-3-31} 所定义的空间 $\XX_{B,P}=\wt{\frakg}\times_{\frakg}\wt{\frakg}_P$ 中的 Zariski 闭包 $\ol X_w$。此时，$G/\mathrm{Res}_{F/\Q_p}B$ 与 $G/\mathrm{Res}_{F/\Q_p}P$ 上两个泛旗之交，并不诱导分次片（相对于 $P$-滤过）上一族良定义的滤过，而只在它们的最高外幂上诱导出良定义的滤过。由此得到图
\[
\begin{tikzcd}[row sep=large,column sep=large]
 & \ol X_w/G\arrow[dl]\arrow[dr] & \\
 \wt{\frakg}/G & & \mathrm{Lie}\,\mathrm{Res}_{F/\Q_p}S_M/\mathrm{Res}_{F/\Q_p}S_M.
\end{tikzcd}
\]

\ysrc{240}
注意 $\ol X_w$ 的构造对包含 $P'\subseteq P$ 具有函子性，这可以从下述等同看出：
\[
X_{B,P}/G=q^{-1}(\frakp)/\mathrm{Res}_{F/\Q_p}P,
\]
其中 $q:\wt{\frakg}\to\frakg$ 是典范投射，$\frakp\subseteq\frakg$ 是 $\mathrm{Res}_{F/\Q_p}P\subseteq G$ 的 Lie 代数。叠 $\ol X_w/G$ 典范同构于 $\ol X'_w/\mathrm{Res}_{F/\Q_p}P$，其中 $\ol X'_w\subseteq q^{-1}(\frakp)$ 是局部闭子集 $V'_w$ 的闭包，而 $V'_w$ 满足 $V_w/G=V'_w/\mathrm{Res}_{F/\Q_p}P$。

注意对固定的 $w$，$G/\mathrm{Res}_{F/\Q_p}B$ 上的泛旗与（$P$ 型）部分标准旗之交，给出标准旗的分次片上一个良定义的旗。这诱导映射 $V'_w\to\wt{\frakm}$，于是我们可以定义 $\wt X_w$ 为 $V'_w$ 在 $q^{-1}(\frakp)\times\wt{\frakm}$ 中的闭包。由此得到图
\[
\begin{tikzcd}[row sep=large,column sep=large]
 & \wt X_w/\mathrm{Res}_{F/\Q_p}P\arrow[dl]\arrow[dr] & \\
 \wt{\frakg}/G & & \wt{\frakm}/\mathrm{Res}_{F/\Q_p}M,
\end{tikzcd}
\]
以及典范投射
\[
\pi_P:\wt X_w/\mathrm{Res}_{F/\Q_p}P\longrightarrow\ol X_w/G,
\]
它与到 $\wt{\frakg}$ 的投射相容。

\begin{prop}{5.3.34}\label{thm:a4-5-3-34}
\begin{enumerate}[label=(\roman*),leftmargin=2.2em,itemsep=0.1em]
\item 若 $P=B$，则概形 $\ol X_w=\wt X_w$ 正规且 Cohen--Macaulay。
\item 对一般的抛物子群 $P$，概形 $\ol X_w$ 单分支（unibranched）。
\end{enumerate}
\end{prop}

关于 $\ol X_w$ 局部几何的这些结果，分别证明于 [BHS19, Prop.\ 2.3.3, Thm.\ 2.3.6] 与 [Wu24, Thm.\ 2.12]。

\begin{hyp}{5.3.35}\label{thm:a4-5-3-35}\tn{底本此处以「Expectation（预期）」为标题，与「Hypothesis（假设）」并列；本译本未另设「预期」环境，故借现有的「假设」环境排版，编号照原，读者应按底本的「预期」理解（上下文推断，无外部依据）。}
\begin{enumerate}[label=(\roman*),leftmargin=2.2em,itemsep=0.1em]
\item 概形 $\ol X_w$ 与 $\wt X_w$ 正规且 Cohen--Macaulay。
\item 态射
\[
\pi_P:\wt X_w/\mathrm{Res}_{F/\Q_p}P\longrightarrow\ol X_w/G
\]
满足 $R\pi_{P,*}\OO=\OO$ 与 $R\pi_{P,*}\omega=\omega$。
\end{enumerate}
\end{hyp}

我们可以用这些对象定义完备化 $\XX_P^{\mathrm{dR},+}$ 沿固定权 $w\lambda$ 的紧化，方式是令
\[
(\ol\XX_P^{\mathrm{dR},+})^{\wedge}_{w\lambda}=(\ol X_w/G)^{\wedge}_0,\qquad
(\wt\XX_P^{\mathrm{dR},+})^{\wedge}_{w\lambda}=(\wt X_w/P)^{\wedge}_0.
\]

\begin{rem}{5.3.36}\label{thm:a4-5-3-36}
当然，正如记号所暗示，我们希望在完备化之前就能定义叠 $\ol\XX_P^{\mathrm{dR},+}$ 与 $\wt\XX_P^{\mathrm{dR},+}$。
\end{rem}

\ysrc{241}
由构造，这些紧化构成下图
\[
\begin{tikzcd}[row sep=large,column sep=1.4em]
 & & (\wt\XX_P^{\mathrm{dR},+})^{\wedge}_{w\lambda}\arrow[ddll]\arrow[dl]\arrow[dr] & \\
 & (\ol\XX_P^{\mathrm{dR},+})^{\wedge}_{w\lambda}\arrow[dl]\arrow[drr] & & (\XX_M^{\mathrm{dR},+})^{\wedge}_{w\lambda}\arrow[d]\\
 (\XX_d^{\mathrm{dR},+})^{\wedge}_{\lambda} & & & (\XX_{S_M}^{\mathrm{dR},+})^{\wedge}_{\det(w\lambda)},
\end{tikzcd}
\]
其中记 $\det(w\lambda)$ 为 $S_M$ 的权，它由 Levi $M$ 的每个 $\mathrm{GL}_m$-块上的最高外幂给出。

最后，下述结果把上述这些叠及其显式局部描述（相应地，其定义）与 $\XX_P$ 的 Drinfeld 紧化重新联系起来。定理的第 (ii) 部分是 [BHS19] 中三角化簇局部模型的推广。

\begin{thm}{5.3.37}\label{thm:a4-5-3-37}
\begin{enumerate}[label=(\roman*),leftmargin=2.2em,itemsep=0.2em]
\item 由 $\XX_P\to\XX_P^{\mathrm{dR},+}$ 在完备化上诱导的典范映射
\[
(\XX_P)^{\wedge}_{w\lambda}\longrightarrow(\XX_P^{\mathrm{dR},+})^{\wedge}_{w\lambda}
\]
延拓为态射
\[
(\ol\XX_P)^{\wedge}_{w\lambda}\longrightarrow(\ol\XX_P^{\mathrm{dR},+})^{\wedge}_{w\lambda}
\]
与
\[
(\wt\XX_P)^{\wedge}_{w\lambda}\longrightarrow(\wt\XX_P^{\mathrm{dR},+})^{\wedge}_{w\lambda}.
\]
\item 设 $\tilde x\in\wt\XX_P$，并记 $\ol x$ 为它在 $\ol\XX_P$ 中的像。假设 $\tilde x$ 在 $\XX_M$ 中的像 $x_M=(D_1,\dots,D_r)$ 是正则的，意思是：
\begin{itemize}[label=--,leftmargin=2.2em,itemsep=0.1em]
\item 映射 $\XX_M\to\XX_M^{\mathrm{dR},+}$ 在 $x_M$ 处形式光滑；
\item 对 $i\neq j$，Ext-群 $\Ext^0(D_i,D_j)=0=\Ext^2(D_i,D_j)$ 消没。
\end{itemize}
则态射
\[
(\ol\XX_P)^{\wedge}_{w\lambda}\longrightarrow(\ol\XX_P^{\mathrm{dR},+})^{\wedge}_{w\lambda}
\]
在 $\ol x$ 处形式光滑。
\end{enumerate}
\end{thm}

\begin{rem}{5.3.38}\label{thm:a4-5-3-38}
在定理的 (i) 情形下，有一个明显的交换图，我们预期它是笛卡尔的。
\end{rem}

\begin{conj}{5.3.39}\label{thm:a4-5-3-39}
在定理的 (ii) 情形下，态射
\[
(\wt\XX_P)^{\wedge}_{w\lambda}\longrightarrow(\wt\XX_P^{\mathrm{dR},+})^{\wedge}_{w\lambda}
\]
在 $\tilde x$ 处形式光滑。
\end{conj}

当然，与关于簇 $\ol X_w$ 局部结构的结果结合起来，这就给出了 Drinfeld 紧化局部几何的一个显式描述。

在 $P=B$ 的情形，下述定理使状况更精确一点：设 $\varphi_1,\dots,\varphi_d\in L^\times$ 满足 $i\neq j$ 时 $\varphi_i/\varphi_j\neq1,q$。记 $\delta_{\mathrm{sm}}$ 为未分歧特征
\[
(t_1,\dots,t_d)\longmapsto\mathrm{unr}_{\varphi_1}(t_1)\cdots\mathrm{unr}_{\varphi_d}(t_d).
\]
对 $w\in W$，我们令
\[
\ol\XX_{B,zw^\lambda\delta_{\mathrm{sm}}}\subseteq\ol\XX_B
\]

\ysrc{242}
为 $z^{w^\lambda}\delta_{\mathrm{sm}}\in\XX_T$ 的原像，并记 $\wh{\ol\XX}_{B,zw^\lambda\delta_{\mathrm{sm}}}\subseteq\ol\XX_B$ 为 $\ol\XX_B$ 沿该闭子叠的形式完备化。我们令
\[
\XX_{d,(\lambda,\delta_{\mathrm{sm}})-\mathrm{tri}}=\bigcup_{w\in W}\ol\beta_B\bigl(\wh{\ol\XX}_{B,zw^\lambda\delta_{\mathrm{sm}}}\bigr)\subseteq\XX_d
\]
为 $\wh{\ol\XX}_{B,zw^\lambda\delta_{\mathrm{sm}}}$ 在 $\XX_d$ 中诸像的并。注意：关于 $\delta_{\mathrm{sm}}$ 的假设蕴涵 $\ol\beta_B$ 在（形式完备化）$\wh{\ol\XX}_{B,zw^\lambda\delta_{\mathrm{sm}}}$ 上的限制是闭嵌入。

按定义，这是 $\XX_d$ 的闭子叠，由所有这样的 $(\varphi,\Gamma)$-模 $D$ 组成：$D$ 带滤过 $\mathrm{Fil}_\bullet$，其分次片为
\[
\mathcal{R}_F\bigl(\mathrm{unr}_{\varphi_i}z(\lambda_{w_\tau(i)})_\tau\bigr)
\]
（记号见 (5.3.18)），其中 $w\in W$ 取某个值。在闭子叠 $\XX_{d,(\lambda,\delta_{\mathrm{sm}})-\mathrm{tri}}$ 上，该滤过并不粘成泛 $(\varphi,\Gamma)$-模 $D$ 的滤过（因为 Weyl 群元素 $w$ 可能变化）。但把 $t$ 取逆之后这就成立，我们记 $\mathcal{M}_\bullet\subseteq\wt D[1/t]$ 为该滤过。有了这个记号，就可以考虑形式子叠
\[
\wh{\XX}_{d,(\lambda,\delta_{\mathrm{sm}})-\mathrm{tri}}\subset\XX_d,
\]
它参数化 $D$ 到 $\XX_{d,(\lambda,\delta_{\mathrm{sm}})-\mathrm{tri}}$ 的无穷小邻域的形变 $\wt D'$，使得 $\mathcal{M}_\bullet$ 提升为（由于我们对 $\delta_{\mathrm{sm}}$ 的选取：必然唯一的）滤过 $\mathcal{M}'_\bullet\subseteq\wt D'[1/t]$，该滤过在 $\varphi$ 与 $\Gamma$ 下稳定。下述定理基本上证明于 [BHS19]。

\begin{thm}{5.3.40}\label{thm:a4-5-3-40}
\begin{enumerate}[label=(\roman*),leftmargin=2.2em,itemsep=0.2em]
\item 态射
\[
\wh{\ol\XX}_{B,zw^\lambda\delta_{\mathrm{sm}}}\longrightarrow(\ol\XX_B^{\mathrm{dR},+})^{\wedge}_{w\lambda}\cong(\ol X_{ww_0}/G)^{\wedge}_0
\]
形式光滑，相对对偶层为 $\mathcal{L}_{2\rho}^{\otimes[F:\Q_p]}$。
\item 存在形式光滑态射
\[
f_{\lambda,\delta_{\mathrm{sm}}}:\wh{\XX}_{d,(\lambda,\delta_{\mathrm{sm}})-\mathrm{tri}}\longrightarrow(\wt{\frakg}\times_{\frakg}\wt{\frakg})^{\wedge}_0,
\]
其相对对偶层为 $\mathcal{L}_{2\rho}^{\otimes[F:\Q_p]}$，且它在 $\ol\beta_B\bigl(\wh{\ol\XX}_{B,zw^\lambda\delta_{\mathrm{sm}}}\bigr)$ 上的限制由 (i) 中的态射给出。
\end{enumerate}
\end{thm}

我们来定义定理陈述中出现的层 $\mathcal{L}_{2\rho}$：叠 $\wh{\ol\XX}_{B,zw^\lambda\delta_{\mathrm{sm}}}$ 有典范映射到 $*/\GG_m^d$（因为它位于 $\ol\XX_B$ 内，而 $\ol\XX_B$ 映到 $\XX_T$，后者又进一步投射到 $*/\GG_m^d$）。因此，拉回正根之和 $2\rho$ 所给出的特征标对应的线丛是有意义的。这在 (i) 中定义了该线丛。为理解 (ii) 中的线丛，我们注意可以把 $\mathcal{L}_{2\rho}$ 写成 $\ol X_{ww_0}$ 上某个 $G$-等变线丛的拉回；后者定义在整个 $\wt{\frakg}\times_{\frakg}\wt{\frakg}\subset\frakg\times G/B\times G/B$ 上，由其中一个 $G/B$-因子上的线丛（对应于 $2\rho$ 的线丛）拉回给出。

\ysrc{243}
\begin{exam}{5.3.41}\label{thm:a4-5-3-41}
我们在 $\mathrm{GL}_2(\Q_p)$ 的情形把上述结果再展开一些。

(i) 设 $\delta_1=\mathrm{unr}_{\varphi_1}z^{\lambda_1}$ 与 $\delta_2=\mathrm{unr}_{\varphi_2}z^{\lambda_2}$ 是两个局部代数特征标，其未分歧光滑部分满足 $\varphi_1/\varphi_2\notin\{1,p^{\pm1}\}$，且 $\lambda_1\neq\lambda_2$。考虑一个 $(\varphi,\Gamma)$-模 $D$，它是扩张
\begin{equation}\tag{5.3.42}\label{eq:a4-5-3-42}
0\longrightarrow\mathcal{R}_L(\delta_1)\longrightarrow D\longrightarrow\mathcal{R}_L(\delta_2)\longrightarrow0.
\end{equation}
基本上有两种情形需要考虑：或者 $D$ 是晶体的，或者扩张非分裂且 $\lambda_1<\lambda_2$。在晶体情形，我们还需要区分 $\lambda_1<\lambda_2$ 与 $\lambda_1>\lambda_2$。
\begin{enumerate}[label=(\alph*),leftmargin=2.2em,itemsep=0.15em]
\item 扩张非分裂且 $\lambda_1<\lambda_2$。此时 $D$ 是 Hodge--Tate 但不是 de Rham，并且 \eqref{eq:a4-5-3-42} 是把 $D$ 写成两个秩一 $(\varphi,\Gamma)$-模的扩张的唯一方式。特别地，$(\delta_1,\delta_2)$ 是唯一一对满足 $(D,\delta_1,\delta_2)\in\ol\XX_B$ 的特征标。但若令 $\delta'_1=\mathrm{unr}_{\varphi_1}z^{\lambda_2}$ 与 $\delta'_2=\mathrm{unr}_{\varphi_2}z^{\lambda_1}$，则 $(D,\delta'_1,\delta'_2)\in\ol\XX_B$。于是
\[
x_1=(D,\delta_1,\delta_2)\qquad\text{与}\qquad x_s=(D,\delta_1,\delta_2)
\]
是 $\ol\XX_B$ 中位于 $D$ 上方的仅有的点\tn{底本此处 $x_s$ 与 $x_1$ 印成同一式 $(D,\delta_1,\delta_2)$，与紧随其后的「$x_1$ 的邻域、$x_s$ 的邻域」以及上文给出的 $\delta'_1,\delta'_2$ 不符，疑为 $(D,\delta'_1,\delta'_2)$ 之误；此处照录底本，并于译文中保留疑点（依据：同页上文与下文语境）。}。因此，在 $D$ 处局部地，空间 $\XX_{d,(\lambda,\delta_{\mathrm{sm}})-\mathrm{tri}}$ 有两个连通分支：一个是 $x_1$ 的（在 $\ol\XX_B$ 内的）邻域在投射 $\ol\XX_B\to\XX_d$ 下的像，另一个是 $x_s$ 的（在 $\ol\XX_B$ 内的）邻域的像。
\item 扩张 $D$ 是晶体的。此时还要再分两种情形：
\begin{enumerate}[label=({b}\arabic*),leftmargin=2.8em,itemsep=0.15em]
\item $\lambda_1>\lambda_2$，且扩张非分裂。此时扩张必然是晶体的。此外，空间 $\XX_{d,(\lambda,\delta_{\mathrm{sm}})-\mathrm{tri}}$ 在 $D$ 处光滑。

$\ol\XX_B$ 中位于 $D$ 上方恰有两个点，且两者实际上都属于 $\XX_B$。第二个点的出现方式与 (a) 中不同：它来自下述事实——$D$ 也可以写成扩张
\[
0\longrightarrow\mathcal{R}_L(\eta_1)\longrightarrow D\longrightarrow\mathcal{R}_L(\eta_2)\longrightarrow0,
\]
其中 $\eta_1=\mathrm{unr}_{\varphi_2}z^{\lambda_1}$，$\eta_2=\mathrm{unr}_{\varphi_1}z^{\lambda_2}$。也就是说：我们改变了加细（refinement），即改变了 $\varphi_1$ 与 $\varphi_2$ 的次序（而在 (a) 中改变的是权 $\lambda_1,\lambda_2$ 的次序）；换在定理 5.3.40 的语言中，这些点会通过另一个光滑特征标 $\delta_{\mathrm{sm}}$ 的选取来处理。
\item $\lambda_1>\lambda_2$ 且扩张分裂；或者 $\lambda_1<\lambda_2$。此时 $\ol\XX_B$ 中有三个点，其中两个对应于（按 $(\varphi_1,\varphi_2)$ 的）次序，即属于 $\XX_{d,(\lambda,\delta_{\mathrm{sm}})-\mathrm{tri}}$。它们的描述与 (a) 中完全一样。(a) 与 (b2) 的唯一差别在于：$\ol\XX_B$ 中还有第三个点映到 $D\in\XX_2$，它对应于（按 $(\varphi_2,\varphi_1)$ 的）次序给出的加细。当然，这个点不被定理 5.3.40 覆盖。
\end{enumerate}
\end{enumerate}

(ii) 我们描述 $d=2$ 情形的局部模型。定理 5.3.40 的局部模型 $\wt{\frakg}\times_{\frakg}\wt{\frakg}$ 有两个不可约分支，二者都非奇异。为看出这一点，也为得到局部模型的显式方程，我们注意 $\wt{\frakg}\times_{\frakg}\wt{\frakg}$ 光滑等价于 $q^{-1}(\frakb)$，其中 $q:\wt{\frakg}\to\frakg$ 是典范投射。
\end{exam}

\ysrc{244}
更确切地说，作为叠有
\[
(\wt{\frakg}\times_{\frakg}\wt{\frakg})/G=(q^{-1}(\frakb))/B.
\]
显然
\[
X_1=\frakb\times\{[1:0]\}\subset q^{-1}(\frakb)\subset\frakb\times\PP^1
\]
是不可约分支之一。另一个分支 $X_s$ 是
\[
V_s=\left\{\left(\begin{pmatrix}a&b\\0&c\end{pmatrix},[x:1]\right)\ \middle|\ (a-c)x+b=0\right\}
\]
的闭包，其中 $s$ 记 $\mathrm{GL}_2$ 的 Weyl 群 $\mathcal{S}_2$ 的非平凡元。注意 $V_s$ 光滑、在 $X_s$ 中开，且与 $X_1$ 不交。交 $X_1\cap X_s$ 含于原点 $0=[1:0]\in\PP^1$ 的原像中，而该点的标准仿射邻域 $q^{-1}(\frakb)$ 由
\[
\left\{\left(\begin{pmatrix}a&b\\0&c\end{pmatrix},[1:y]\right)\ \middle|\ y\bigl((a-c)+by\bigr)=0\right\}
\]
描述，而 $X_s$ 由方程
\[
\left\{\left(\begin{pmatrix}a&b\\0&c\end{pmatrix},[1:y]\right)\ \middle|\ (a-c)+by=0\right\}
\]
给出。

我们记
\[
Z=X_1\cap X_s=\left\{\left(\begin{pmatrix}a&b\\0&a\end{pmatrix},[1:0]\right)\right\}
\]
为两个分支之交。在定理 5.3.40 的形式光滑态射下，本例第一部分中类型 (a) 与 (b2) 的点映到 $Z$，而类型 (b1) 的点映到
\[
\left\{\left(\begin{pmatrix}a&0\\0&a\end{pmatrix},[x:1]\right)\right\}\subset V_s.
\]
更确切地说，类型 (a) 的点映到矩阵
\[
\begin{pmatrix}a&b\\0&a\end{pmatrix}
\]
（其中 $b\neq0$），而类型 (b2) 的点映到对角矩阵 $\mathrm{diag}(a,a)$。

(iii) $\ol\XX_B$ 在 $\XX_d\times\XX_T$ 中的像 $\XX_{d,\mathrm{tri}}$ 应当看成特征簇的局部伽罗瓦论对应物。关于这一口号的数学上有意义的表述，参看 [BHS17b]。在特征曲线（eigencurve）的情形，我们简述 (i) 中讨论的点如何出现。设 $f$ 是权 $k\ge2$、级 $\Gamma_1(N)$ 的经典模特征形式，其中 $p\nmid N$。则 $f$ 所附的伽罗瓦表示 $\rho$ 在 $p$ 处是晶体的。在级 $\Gamma_1(N)\cap\Gamma_0(p)$ 处，$f$ 有两个稳定化 $f_\alpha$ 与 $f_\beta$；它们是 $p$ 之外同一组 Hecke 特征值的特征形式，但 $U_p$-特征值不同，记为 $\alpha$ 与 $\beta$。这两个不同的 $U_p$-特征值对应于上述加细的两种可能选取。在伽罗瓦表示形变空间的一般纤维中，特征曲线至少有两条分支在 $\rho$ 处相遇：含经典模形式 $f_\alpha$ 的分支，以及含经典形式 $f_\beta$ 的分支。

当 $\rho$ 在 $p$ 处非分裂时，这就是含伽罗瓦表示 $\rho$ 的全部支。这对应于上面情形 (b1) 中的非分裂扩张（注意 $f_\alpha$ 与 $f_\beta$ 的权都是支配代数权）。但若 $\rho$ 在 $p$ 处分裂，\ysrc{245}
则在 $p$ 处晶体时，特征曲线还有第三支，含非经典 $p$-进模形式 $f'$（权为 $1-k$）。这正是上面讨论的情形 (b2)。我们明确指出：定理 5.3.40 中在给定 $(\varphi,\Gamma)$-模 $D$ 处相遇的那两条 $\wh{\XX}_{d,(\lambda,\delta_{\mathrm{sm}})-\mathrm{tri}}$ 的分支（如 (b2) 所述），并不是含经典形式 $f_\alpha$ 与 $f_\beta$ 的分支——应把其中一支看成含经典模形式，另一支含非经典形式 $f'$（该分支上的点具有非经典权）。最后我们指出：(a) 中的点对应于非经典 $p$-进模形式（它们可能有也可能没有经典权）。

\begin{rem}{5.3.43}\label{thm:a4-5-3-43}
在 [BD24] 中，Breuil 与 Ding 定义了三角化簇的一个变体，其对象是 paraboline $(\varphi,\Gamma)$-模\tn{「paraboline」为 Breuil–Ding 的专门术语，未见于通行中译，此处保留原文（上下文推断，无外部依据）。}。在上述叠的语境中，他们的构造可以描述如下（记号与 Breuil--Ding 的不同）：我们把 $\XX_M$ 换成它的闭子叠 $\XX'_M$，后者的对象是 de Rham、具有给定权与惯性型的对象，模去与任意秩一对象的扭。于是可以在 $\XX_P$ 中定义 $\XX'_M$ 的原像 $\XX'_P$，并取其闭包 $\ol\XX'_P$ 于紧化 $\ol\XX_P$ 中。那么在本文的语言中，Breuil--Ding 的 paraboline 簇的类比就是 $\ol\XX'_P$ 在 $\XX_d\times\XX_{S_M}$ 中的像。Breuil 与 Ding 还为他们的空间构造了局部模型。这些局部模型由下述闭子空间在（\eqref{eq:a4-5-3-31} 的空间 $X_{B,P}$ 内的）Zariski 闭包给出：$V_w$ 中像落在 $\frakm$ 的中心的那些元素。人们或许会问：$\wt\XX_P$ 是否有类似的变体？这一变体未出现在 Breuil--Ding 的工作中，但 Huang [Hua23] 的工作中对此有讨论。
\end{rem}

\section{范畴 $p$-进局部 Langlands 猜想}
本节陈述我们的主要猜想，并解释其部分动机（特别是一些来自 Taylor--Wiles 方法的动机）。

\subsection{预期：Banach 情形}
我们现在把 Banach 情形中的预期说清楚。为此先引进一些记号。与通常一样，固定有限扩张 $F/\Q_p$ 与系数环 $\OO$，令 $G=\mathrm{GL}_d(F)$，$K=\mathrm{GL}_d(\OO_F)$，$Z=Z(G)$。

如我们已经解释过的，目标是建立一个猜想，把（光滑表示的导出范畴与）凝聚层、更一般地 $(\XX_d$ 这一形式代数叠上的）Ind-凝聚复形联系起来。为把陈述说精确，我们需要先讨论 $G$-表示的相关范畴。

\subsubsection{表示范畴的预备}
我们从回忆附录 E 的部分材料开始。以 $\mathrm{sm}\,G$ 记 $G$ 在 $\OO$-模上的光滑表示构成的阿贝尔范畴（即光滑表示范畴），其中 $\OO$-线性 $G$-表示称为光滑的，如果表示中每个元素都被 $G$ 的某个开子群固定，并被 $p$ 的某个幂零化。（见定义 E.1.2。）以 $D(\mathrm{sm}\,G)$ 记 $\mathrm{sm}\,G$ 的导出稳定 $\infty$-范畴。

任一光滑 $G$-表示带有典范的 $\OO[[G]]$-模结构，其中 $\OO[[G]]$ 是定义 E.1.1 中所定义的环；而经常有用的是\ysrc{246}从 $\OO[[G]]$-模结构的角度研究光滑 $G$-表示。为此，命题 E.2.2 给出一个全忠实 t-正合函子
\begin{equation}\tag{6.1.2}\label{eq:a4-6-1-2}
D(\mathrm{sm}\,G)\hookrightarrow D(\OO[[G]])
\end{equation}
（目标是 $\OO[[G]]$-模复形的稳定 $\infty$-范畴），其本质像等于 $D_{\mathrm{sm}}(\OO[[G]])$（即由上同调光滑的对象构成的全子范畴）。

由于存在许多不光滑的 $\OO[[G]]$-模（例如 $\OO[[G]]$ 自身），函子 \eqref{eq:a4-6-1-2} 远不是等价。另一方面，借助对 pro-对象的考虑，可以构造一个函子，它接近于是 \eqref{eq:a4-6-1-2} 的左逆。事实上，若以 $D^-_{\mathrm{c.g.}}(\OO[[G]])$ 记 $D(\OO[[G]])$ 的全子范畴，其对象的上同调作为 $\OO[[G]]$-模可数生成、且在充分高的次数上消没，则 (E.4.4) 给出函子
\begin{equation}\tag{6.1.3}\label{eq:a4-6-1-3}
D^-_{\mathrm{c.g.}}(\OO[[G]])\longrightarrow\mathrm{Pro}\,D^b(\mathrm{sm}\,G).
\end{equation}
（例如，$\OO[[G]]$ 映到 $\mathrm{Pro}\,\mathrm{sm}\,G$ 的形式 pro 对象「$\varprojlim_{H,n}"\,c\text{-}\mathrm{Ind}_H^G\OO/p^n$；这里 $H$ 跑遍 $G$ 的所有紧开子群，$n$ 跑遍所有正整数，而 $\OO/p^n$ 带平凡 $H$-作用。）此外，复合
\[
\begin{aligned}
D^b_{\mathrm{c.g.}}(\mathrm{sm}\,G)&\xrightarrow[(E.4.5)]{\ \sim\ }D^b_{\mathrm{c.g.}\,\mathrm{sm}}(\OO[[G]])\\
&\hookrightarrow D^-_{\mathrm{c.g.}}(\OO[[G]])\xrightarrow{\eqref{eq:a4-6-1-3}}\mathrm{Pro}\,D^b(\mathrm{sm}\,G)
\end{aligned}
\]
恰是典范全忠实函子 $D^b_{\mathrm{c.g.}}(\mathrm{sm}\,G)\hookrightarrow\mathrm{Pro}\,D^b(\mathrm{sm}\,G)$。（这说明了我们在什么意义上认为 \eqref{eq:a4-6-1-3} 接近于 \eqref{eq:a4-6-1-2} 的左逆。）

为正确陈述我们的范畴 Langlands 猜想，我们先给出一个关于阿贝尔范畴 $\mathrm{sm}\,G$ 的预备猜想。

\begin{conj}{6.1.4}\label{thm:a4-6-1-4}
光滑 $\OO$-线性 $G$-表示构成的阿贝尔范畴 $\mathrm{sm}\,G$ 局部凝聚（locally coherent）。
\end{conj}

\begin{rem}{6.1.5}\label{thm:a4-6-1-5}
回忆这意味着 $\mathrm{sm}\,G$ 紧生成，且紧对象构成 $\mathrm{sm}\,G$ 的一个阿贝尔子范畴。我们在命题 E.3.5 中已看到，$\mathrm{sm}\,G$ 中的紧对象恰是有限表现的表示（定义 E.3.3 意义下），且 $\mathrm{sm}\,G$ 由这些对象生成。因此猜想 6.1.4 的内容是：$\mathrm{sm}\,G$ 中由有限表现表示构成的全子范畴对取核封闭。（它显然对取余核封闭。）
\end{rem}

\begin{rem}{6.1.6}\label{thm:a4-6-1-6}
由于有限表现的光滑表示恰是作为 $\OO[[G]]$-模有限表现的表示（引理 E.3.4），若范畴 $\OO[[G]]\text{-}\Mod$（左 $\OO[[G]]$-模）局部凝聚，即若 $\OO[[G]]$ 是凝聚环，也就是有限表现 $\OO[[G]]$-模构成的范畴对取核稳定，则猜想 6.1.4 当然成立。

在 $d=1$ 情形，环 $\OO[[G]]$ 甚至是诺特的；而 $d=2$ 情形由 Timmins 的结果 [Tim23, Thm.\ 9.7] 知其凝聚，该结果基于 Shotton 关于环 $\OO[[\mathrm{SL}_2(F)]]$ 凝聚性的证明 [Sho20, Cor.\ 4.4]。

无论哪种情形，由此推出有限表现光滑 $\OO[[G]]$-模构成的范畴是阿贝尔范畴。

猜想 6.1.4 没有其他已知情形。已知当 $d>2$ 时 $\OO[[G]]$ 不凝聚，这是 Timmins 的结果 [Tim23]。我们相信猜想 6.1.4 的主要理由是\ysrc{247}
我们相信类似猜想 6.1.15 的东西必定成立，而若猜想 6.1.4 不成立，就很难一致地陈述这样的猜想。
\end{rem}

现在设猜想 6.1.4 成立，并以 $D^b_{\mathrm{f.p.}}(\mathrm{sm}\,G)$ 记 $D(\mathrm{sm}\,G)$ 的全子范畴，它由（上同调）有界、且上同调为有限表现 $\OO[[G]]$-模的复形组成。

正如命题 E.3.5 之后的讨论所指出的，此时 $D^b_{\mathrm{f.p.}}(\mathrm{sm}\,G)$ 恰是 $D(\mathrm{sm}\,G)$ 中由凝聚对象构成的全子范畴。

对 $D^b_{\mathrm{f.p.}}(\mathrm{sm}\,G)$ 有对偶理论，定义如下。

我们通过 $g\mapsto g^{-1}$ 把 $\OO[[G]]^{\mathrm{op}}$ 与 $\OO[[G]]$ 等同。这样就能把右 $\OO[[G]]$-模典范地转成左 $G$-模。于是我们定义反等价
\[
\mathbf{D}:D^b_{\mathrm{f.p.}}(\mathrm{sm}\,G)\longrightarrow D^b_{\mathrm{f.p.}}(\mathrm{sm}\,G)
\]
通过
\[
\mathbf{D}(-):=R\Hom_{\OO[[G]]}\bigl(-,\OO[[G]]\bigr)[d^2[F:\Q_p]+1].
\]
我们预期它有幅度 $[0,d]$，带有自然同构 $\mathbf{D}\circ\mathbf{D}\isom\mathrm{Id}$，并具有下述性质：
\begin{equation}\tag{6.1.7}\label{eq:a4-6-1-7}
\mathbf{D}\bigl(c\text{-}\mathrm{Ind}_H^G V\bigr)=c\text{-}\mathrm{Ind}_H^G V^\vee
\end{equation}
对 $G$ 的紧开子群 $H$ 与任意有限长光滑表示 $V$ 成立；这里 $V^\vee=\Hom_\OO(V,E/\OO)$ 记 $V$ 的庞特里亚金对偶，带有其反变 $H$-作用。

\subsubsection{$\XX_d$ 上的 Serre 对偶}
在伽罗瓦侧，我们有（Ind-）凝聚层的导出范畴 $D_{\mathrm{coh}}(\XX_d)$ 与 $\IndCoh(\XX_d)$（见附录 B.2）。$\XX_d$ 上有一个对合 $\imath$，在模论意义下由 $\rho\mapsto\rho^\vee$ 给出。以 $D_{\XX_d}$ 记 $D_{\mathrm{coh}}(\XX_d)$ 的反等价，它由 Grothendieck--Serre 对偶与 $\imath^*$ 复合而成，即
\[
\imath^*\circ R\mathcal{H}om_{\OO_{\XX_d}}(-,\omega_{\XX_d}),
\]
其中 $\omega_{\XX_d}$ 是 $\XX_d$ 的对偶层，看作 $\XX$ 上的 pro-凝聚层（如例 B.3.2 所讨论），放在 $0$ 次；回忆我们预期 $\XX_d$ 是 lci（局部完全交）。

\subsubsection{Hodge 型、惯性型、Bernstein 中心的连通分支与 Hecke 代数}
由 Schneider--Zink [SZ99] 的结果（「惯性局部 Langlands 对应」），可以对每个惯性型 $\tau:I_F\to\mathrm{GL}_d(L)$ 附以一个有限维光滑不可约 $K$-表示 $\sigma^{\mathrm{crys}}(\tau)$，见 [EG23, Thm.\ 8.2.1]。以 $\sigma^{\mathrm{crys},\circ}(\tau)$ 记 $\sigma^{\mathrm{crys}}(\tau)$ 中某个 $K$-稳定的 $\OO$-格。如 [CEG+16] 中所解释，$\tau$ 决定一个 Bernstein 分支 $\Omega_\tau$，且有典范同构
\begin{equation}\tag{6.1.10}\label{eq:a4-6-1-10}
Z_{\Omega_\tau}\isom\End_G\bigl(c\text{-}\mathrm{Ind}_K^G\sigma^{\mathrm{crys}}(\tau)\bigr).
\end{equation}
（这里及下文，对任意 $K$-表示 $U$，$c\text{-}\mathrm{Ind}_K^G U$ 取通常含义：它是紧支撑函数 $G\to U$ 构成的空间，这些函数关于 $G$ 在左侧的 $K$-作用是 $K$-等变的，并带由右平移给出的 $G$-作用。）若 $\lambda$ 是正则 Hodge 型，则有 $K$ 的代数表示 $W_\lambda$（见定义 2.3.2）。我们令
\[
\Pi(\lambda,\tau)^\circ:=c\text{-}\mathrm{Ind}_K^G\bigl(W_\lambda\otimes_\OO\sigma^{\mathrm{crys},\circ}(\tau)\bigr)
\]
（注意这依赖于格 $\sigma^{\mathrm{crys}}(\tau)$ 的选取），并令
\[
\Pi(\lambda,\tau):=\Pi(\lambda,\tau)^\circ[1/p].
\]

\ysrc{248}
由于 $W_\lambda[1/p]$ 自然地是 $G$-表示，紧诱导的「推拉」（push-pull）表明
\begin{equation}\tag{6.1.11}\label{eq:a4-6-1-11}
\begin{aligned}
\Pi(\lambda,\tau)&:=\Bigl(c\text{-}\mathrm{Ind}_K^G\bigl(W_\lambda\otimes_\OO\sigma^{\mathrm{crys},\circ}(\tau)\bigr)\Bigr)[1/p]\\
&=c\text{-}\mathrm{Ind}_K^G\bigl(W_\lambda[1/p]\otimes_L\sigma^{\mathrm{crys}}(\tau)\bigr)
=W_\lambda[1/p]\otimes_L c\text{-}\mathrm{Ind}_K^G\sigma^{\mathrm{crys}}(\tau).
\end{aligned}
\end{equation}
由于 $W_\lambda[1/p]$ 还是 $G$ 的不可约代数表示，而紧诱导是光滑 $G$-表示，于是我们得到
\begin{equation}\tag{6.1.12}\label{eq:a4-6-1-12}
\begin{aligned}
\End_G\bigl(\Pi(\lambda,\tau)\bigr)&=\End_G\bigl(W_\lambda[1/p]\otimes_L c\text{-}\mathrm{Ind}_K^G\sigma^{\mathrm{crys}}(\tau)\bigr)\\
&=\End_G\bigl(c\text{-}\mathrm{Ind}_K^G\sigma^{\mathrm{crys}}(\tau)\bigr)\overset{\eqref{eq:a4-6-1-10}}{=}\mathcal{Z}_{\Omega_\tau}.
\end{aligned}
\end{equation}
正如 §5.2.11 中所回忆的，Bernstein 中心 $\mathcal{Z}_{\Omega_\tau}$ 也可以等同于 Weil--Deligne 表示构成的叠上的整体函数，且如在 (5.2.16) 中那样，有态射
\begin{equation}\tag{6.1.13}\label{eq:a4-6-1-13}
\mathcal{Z}_{\Omega_\tau}\longrightarrow\Gamma\bigl((\XX_d^{\mathrm{crys},\lambda,\tau})^{\mathrm{rig}}_\eta,\OO_{(\XX_d^{\mathrm{crys},\lambda,\tau})^{\mathrm{rig}}_\eta}\bigr).
\end{equation}

\subsubsection{猜想的陈述}
我们现在可以陈述范畴 Langlands 猜想。

\begin{conj}{6.1.15}\label{thm:a4-6-1-15}
假设猜想 6.1.4 成立。则存在 $\OO$-线性正合全忠实函子 $\mathfrak{A}:D^b_{\mathrm{f.p.}}(\mathrm{sm}\,G)\to D^b_{\mathrm{coh}}(\XX_d)$，满足下述性质：
\begin{enumerate}[label=(\arabic*),leftmargin=2.2em,itemsep=0.15em]
\item 函子 $\mathfrak{A}$ 有有界的上同调幅度。
\item（$L_\infty$ 是层。）$L_\infty:=\mathfrak{A}(\OO[[G]])$ 是 pro-凝聚层，集中于 $0$ 次。此外它 $\OO[[K]]$-平坦。
\item（与对偶相容。）存在反变函子之间的自然等价
\begin{equation}\tag{6.1.16}\label{eq:a4-6-1-16}
\mathfrak{A}\circ\mathbf{D}\isom(\mathbf{D}_{\XX_d}\circ\mathfrak{A})\bigl[[F:\Q_p]d(d+1)/2+1\bigr],
\end{equation}
（从 $D^b_{\mathrm{f.p.}}(\mathrm{sm}\,G)$ 到 $D_{\mathrm{coh}}(\XX_d)$ 的）。
\item（局部代数向量。）对任意正则 Hodge 型 $\lambda$ 与惯性型 $\tau$，$\wh{\mathfrak{A}\bigl(\Pi(\lambda,\tau)^\circ\bigr)}$ 的概形论支撑等于 $\XX_d^{\mathrm{crys},\lambda,\tau}$。（此处的尖帽记 $p$-进完备化。）此外，由等同 \eqref{eq:a4-6-1-12} 与 $\mathfrak{A}$ 的函子性在 $\mathfrak{A}(\Pi(\lambda,\tau)^\circ)[1/p]$ 上诱导的 Bernstein 中心 $\mathcal{Z}_{\Omega_\tau}$ 的作用，与 \eqref{eq:a4-6-1-13} 给出的作用一致。
\end{enumerate}
\end{conj}

\begin{rem}{6.1.17}\label{thm:a4-6-1-17}
猜想 6.1.15 中函子的记号 $\mathfrak{A}$ 取自论文 [Zhu25a]，那里在 $\ell\neq p$ 情形使用了同样的记号。
\end{rem}

\begin{rem}{6.1.18}\label{thm:a4-6-1-18}
假设 $\mathfrak{A}$ 存在，则它有自然的 Ind 与 Pro 延拓；考虑这些延拓是重要的（实际上也是使所陈述的猜想有意义所必需的）。先看 $\mathfrak{A}$ 的 Ind 延拓；它是连续函子
\[
\mathrm{Ind}\,D^b_{\mathrm{f.p.}}(\mathrm{sm}\,G)\longrightarrow\mathrm{Ind}\,\mathrm{Coh}(\XX),
\]
我们仍记之为 $\mathfrak{A}$\tn{底本此处以 $\XX$（即 $\mathcal{X}$）记 $\XX_d$，与他处记号混用；译文照原保留，后文同类情形不再出注（依据：同页与下页均出现 $\mathrm{Ind}\,\mathrm{Coh}(\XX)$ 与 $\mathrm{Pro}\,\mathrm{Coh}(\XX_d)$ 两种写法）。}。
\end{rem}

\ysrc{249}
范畴 $\mathrm{Ind}\,D^b_{\mathrm{f.p.}}(\mathrm{sm}\,G)$ 在我们的语境中是 [Zhu25a, §4.1] 所考虑的光滑表示重整化范畴的类比。它带有一个典范函子
\begin{equation}\tag{6.1.19}\label{eq:a4-6-1-19}
\mathrm{Ind}\,D^b_{\mathrm{f.p.}}(\mathrm{sm}\,G)\longrightarrow D(\mathrm{sm}\,G),
\end{equation}
若在两侧都限制到下有界复形，它就成为等价。事实上，假设猜想 6.1.4，$D(\mathrm{sm}\,G)$ 上的 t-结构是凝聚的（定义 A.6.7 意义下），而 $D^b_{\mathrm{f.p.}}(\mathrm{sm}\,G)$ 恰是 $D(\mathrm{sm}\,G)$ 中凝聚对象构成的子范畴，于是该函子是 (A.6.8) 的一个特例。特别地，$\mathrm{Ind}\,D^b_{\mathrm{f.p.}}(\mathrm{sm}\,G)$ 的中心与 $D(\mathrm{sm}\,G)$ 的中心一致，即与阿贝尔范畴 $\mathrm{sm}\,G$ 一致。因此，把 $\mathfrak{A}$ 延拓到整个 $\mathrm{Ind}\,D_{\mathrm{f.p.}}(\mathrm{sm}\,G)$，特别使我们能对所有的光滑 $G$-表示计算 $\mathfrak{A}$，而不只是对有限表现的表示。

顺带指出，一般而言 $D(\mathrm{sm}\,G)$ 上的 t-结构不必正则（定义 A.6.9 意义下），因此 \eqref{eq:a4-6-1-19} 一般不必是等价。事实上，我们预期 $D(\mathrm{sm}\,G)$ 上的 t-结构正则、从而 \eqref{eq:a4-6-1-19} 是等价，恰当 $I_1$（$K$ 的 pro-$p$-Iwahori 子群）无挠。

由于 $D^b(\mathrm{sm}\,G)$ 是 $\mathrm{Ind}\,D^b_{\mathrm{f.p.}}(\mathrm{sm}\,G)$ 的全子范畴，我们可以把 $\mathfrak{A}$（的 Ind 延拓）限制到 $D^b(\mathrm{sm}\,G)$，从而得到函子 $D^b(\mathrm{sm}\,G)\to\mathrm{Ind}\,\mathrm{Coh}(\XX)$。（实际上，考虑到注 6.1.28，我们看到该函子取值于 $D^b_{\mathrm{coh}}(\XX)$。）我们可以再把这个函子 Pro 延拓，得到函子
\[
\mathrm{Pro}\,D^b(\mathrm{sm}\,G)\longrightarrow\mathrm{Pro}\,\mathrm{Ind}\,\mathrm{Coh}(\XX).
\]
把它与函子 \eqref{eq:a4-6-1-3} 复合，我们看到对 $D^-_{\mathrm{c.g.}}(\OO[[G]])$ 的对象计算 $\mathfrak{A}$ 是有意义的。§6.1.1 的讨论表明，该函子在 $D^b_{\mathrm{f.p.}}(\mathrm{sm}\,G)\isom D^b_{\mathrm{f.p.}\,\mathrm{sm}}(\OO[[G]])$ 上的限制与原函子 $\mathfrak{A}$ 一致（这里目标取显然的含义，即 $D^b_{\mathrm{sm}}(\OO[[G]])$ 的全子范畴，由上同调有限表现的对象组成；所指的等价直接由命题 E.2.2 推出：把那里陈述的等价限制到其源与目标中上同调有限生成的对象构成的全子范畴即可）。

\begin{rem}{6.1.20}\label{thm:a4-6-1-20}
为使猜想 6.1.15 (2) 有意义，注意如注 6.1.18 所讨论，我们可以对 $\OO[[G]]$（看作自身上的模）计算 $\mathfrak{A}$。先验地，这给出 $\mathrm{Pro}\,\mathrm{Ind}(\XX_d)$ 的一个对象，但 (2) 断言我们实际上得到 $\mathrm{Pro}\,\mathrm{Coh}(\XX_d)$ 的一个对象。
\end{rem}

\begin{rem}{6.1.21}\label{thm:a4-6-1-21}
如 §1.4 已提到的，猜想 6.1.15 看来可以看作某个更强猜想的推论，该更强猜想说：应当存在如 Fargues--Scholze [FS24, Conj.\ I.10.2] 猜想中的（导出）范畴的等价，但系数是 $p$-进而非 $\ell$-进。然而在撰写时似乎还不清楚：为得到这样的等价，究竟应当考虑 $\Bun_G$ 上哪个层范畴。
\end{rem}

\begin{rem}{6.1.22}\label{thm:a4-6-1-22}
$\mathfrak{A}$ 的全忠实性特别蕴涵：$p$-进导出 Hecke 代数与 $\XX_d$ 上某些凝聚层的自同态环之间应当有 $E_1$-环的等同（例如 Feng 的谱 Hecke 代数 [Fen19]）；$\ell\neq p$ 情形的类似猜想见 [Zhu25a, Conj.\ 4.3.1]。
\end{rem}

\ysrc{250}
\begin{rem}{6.1.23}\label{thm:a4-6-1-23}
我们同样可以对潜在半稳定表示陈述猜想 6.1.15 (4)，只需修改表示 $\sigma^{\mathrm{crys}}(\tau)$（例如见 [Pyv21]）；但为叙述简便，我们限于潜在晶体情形。
\end{rem}

\begin{rem}{6.1.24}\label{thm:a4-6-1-24}
猜想 6.1.15 的函子必然是附录 C.2 所考虑的形状，即
\[
\pi\longmapsto L_\infty\otimes^L_{\OO[[G]]}\pi,
\]
其中如在猜想 6.1.15 (2) 中，$L_\infty:=\mathfrak{A}(\OO[[G]])$ 是模叠 $\XX_d$ 上的 pro-凝聚层，并且还是右 $\OO[[G]]$-模（因为 $\OO[[G]]$ 的自同态代数为 $\OO[[G]]$）。（在 $d=1$ 情形，我们将在下文看到 $L_\infty$ 是 $\XX_1$ 上的结构层，见 (7.1.2)；在 $\mathrm{GL}_2(\Q_p)$ 情形我们给出 $L_\infty$ 的构造；但一般情形我们没有 $L_\infty$ 的猜想候选。）

事实上，假设猜想 6.1.4，我们可以用若干个 $\OO[[G]]$ 的直和构成的（可能无限的）复形来消解任一有限表现光滑 $\OO[[G]]$-模 $\pi$
\[
\cdots\longrightarrow\OO[[G]]^{\oplus m_n}\cdots\longrightarrow\OO[[G]]^{\oplus m_0}\longrightarrow\pi\longrightarrow0,
\]
并发现 $\mathfrak{A}(\pi)$ 由
\[
\cdots\longrightarrow\mathfrak{A}(\OO[[G]])^{\oplus m_n}\cdots\longrightarrow\mathfrak{A}(\OO[[G]])^{\oplus m_0}
\]
计算，从而由
\[
\begin{aligned}
&L_\infty\otimes_{\OO[[G]]}\bigl(\cdots\to\OO[[G]]^{\oplus m_n}\cdots\to\OO[[G]]^{\oplus m_0}\bigr)\\
&\qquad=L_\infty\otimes^L_{\OO[[G]]}\pi
\end{aligned}
\]
计算。（假设 (1)，即 $\mathfrak{A}$ 幅度的有界性，为对这些可能下方无界的复形所作的运算提供了依据。）由于 $\mathrm{Ind}\,D^b_{\mathrm{f.p.}}(\mathrm{sm}\,G)$ 在平移与余极限下由有限表现的表示生成，一般地有
\begin{equation}\tag{6.1.25}\label{eq:a4-6-1-25}
\mathfrak{A}(\pi)=L_\infty\otimes^L_{\OO[[G]]}\pi.
\end{equation}
\end{rem}

\begin{rem}{6.1.26}\label{thm:a4-6-1-26}
我们指出：先验地假设 $\mathfrak{A}$ 的幅度有界（即在猜想中强加条件 (1)）的唯一理由，是为使导致公式 \eqref{eq:a4-6-1-25} 的运算有依据；而我们应当强调，对 $\mathfrak{A}$ 存在这样的公式，是我们看待猜想 6.1.15 的根本。若承认该公式成立，则 $\mathfrak{A}$ 幅度的有界性应当随之成立，如下文注 6.1.28 所解释。
\end{rem}

\begin{rem}{6.1.27}\label{thm:a4-6-1-27}
若 $V$ 是拓扑 $\OO[[K]]$-模，且作为 $\OO$-模有限型（因此是 $p$-进完备的），则我们可以把 $p$-进完备的紧诱导表示 $(c\text{-}\mathrm{Ind}_K^GV)^\wedge$ 看作 $\mathrm{Pro}\,\mathrm{sm}\,G$ 的对象，从而构成 $\mathfrak{A}\bigl((c\text{-}\mathrm{Ind}_K^GV)^\wedge\bigr)$。（这也说明了我们在猜想 6.1.15 的性质 (4) 中如何构成 $\mathfrak{A}\bigl(\wh{\Pi(\lambda,\tau)^\circ}\bigr)$。）
\end{rem}
